\section{Political Platonism}

\begin{quotex}
The philosophical background of Guenon and his followers' traditionalism is extremely close to the Platonic tradition ... In other words, traditionalism can be taken as radical Platonism.
\flright{\textsc{Alexander Dugin}, \emph{Political Platonism}}
\end{quotex}


Dugin's purpose seems to be to make Tradition academically palatable, by couching it in Platonic language. Obviously, as we shall demonstrate, unlike Dugin, there are certainly correspondences with Plato. However, Guenon probably refers to Aristotle and the Scholastics more than to Plato, so we shall also rely on them in developing a ``Platonic'' political system. Unfortunately, Dugin makes this rather strange assertion:

\begin{quotex}
Traditionalism offers the entire necessary philosophical, ideational, conceptual, and sociological apparatus to oppose globalization.
\end{quotex}


Au contraire, Guenon rejected philosophy as now understood; a fortiori, he had no intention of creating any such apparatus, since it would be a symptom of modernity, not its cure. That is because Dugin presumes that Guenon created a ``school'' called ``Traditionalism''. That would make his project just one opinion among many. Now Guenon anticipated that; this is what he wrote about ``Traditionalism'':

\begin{quotex} 
``Traditionalists'' refer to people who only have a sort of tendency or aspiration towards tradition \emph{without really knowing anything at all about it}; this is the measure of the distance dividing the ``traditionalist'' spirit from the truly traditional spirit, \emph{for the latter implies a real knowledge}.\footnote{See \textit{The Absurdity of Traditionalism}: \url{https://www.gornahoor.net/?p=728}}
\end{quotex}


I hate to sound harsh, but those are Guenon's words. Dugin is not alone; even Charles Upton refers to a ``School of Traditionalism''. And they are hardly alone in promoting that misunderstanding. Reputations are built on the ignorance of readers.

Since we are on the topic of Platonism, we must understand the distinction he made between opinion and real knowledge. Guenon is not interested in opinion, since opinions can be mistaken. Most thinkers are, since it is more fun to debate. On the other hand, knowledge is of Being or Reality and cannot be mistaken. Guenon calls the transition from opinion to knowledge as an ``intellectual conversion'' or ``metaphysical realization''. For the \emph{soi disant} Traditionalists, such knowledge is simply invisible to them, beyond their intellectual horizon, so they fail to understand it, and, worse, don't even seek it.

\paragraph{Structure of Man and Society}


We begin a Platonic political system with an understanding of the Being of man, who is constituted by the three soul forces, and united by a Self:

\begin{itemize}


\item Concupiscence (eros)

\item Irascibility (thumos)

\item Intellect (nous)

\item The I
\end{itemize}


This is not a matter of opinion, but rather true knowledge of man's Being. When this is applied to political life, we get the idea of castes:

\begin{itemize}


\item Workers and producers

\item Those holding political power

\item Those wielding spiritual authority

\item The (one) leader, chief, king, emperor, etc.
\end{itemize}


\textbf{Georges Dumezil} documented the same socio-political structures in Traditional societies. Hence, those structures are not at all arbitrary, but follow from the Real structure of man and society. Alternative political systems are simply deformations of this. Usually, the political powers and spiritual authorities are not explicitly acknowledged but are illicitly smuggled in.

\paragraph{Being and Thought}


As Thomas Aquinas points out in \emph{De Principiis Naturae}\footnote{\url{https://dhspriory.org/thomas/DePrincNaturae.htm}}, generation, or Being, requires three principles:

\begin{itemize}


\item being in potency which is matter (Prime Matter, or what Dugin calls Chaos)

\item non-existence in act which is privation

\item that through which something comes to be in act which is form (or essences)
\end{itemize}


That is, Existence is generated by essences imposing themselves on prime matter. We experience existence, but know essences. Modern dualistic philosophies separate existence and essence. However, there is no duality, even though the modern mind has trouble discerning essences. Spiritual vision must be developed in order to intuit the essences the generate phenomena.

Privation, which is the negation of form, adds a level of difficulty. Privation is known in Thought, but confused in Being. Therefore, care must be taken to distinguish the part of a being that is the manifestation of the form from privation which is the inability of the form to completely manifest itself. Only in God is there no privation, since His existence and essence are identical.

\paragraph{Struggle for Existence}


Although some still hope for some Utopian time of peace, the struggle for Existence is part of the fabric of nature. Even in Eden, Adam was expected to work. That means he was incomplete, subject to privation, and still needed to manifest all his possibilities.

Moreover, there was temptation in the Garden in the form of the forbidden fruit and the serpent. Therefore, Adam still had to face the Greater Battle against his lower nature.

There is a persistent idea that the animals had a different nature from what we see now. However, the naming of the animals means that Adam understood their true nature. The beasts are wholly natural beings, with no supernatural graces, so the lions were certainly not dining on Impossible Burgers. Thomas Aquinas offers this rationale:

\begin{quotex}
In the opinion of some, those animals which now are fierce and kill others, would, in that state, have been tame, not only in regard to man, but also in regard to other animals. But this is quite unreasonable. For the nature of animals was not changed by man's sin, as if those whose nature now it is to devour the flesh of others, would then have lived on herbs, as the lion and falcon.
\flright{\textit{Summa Theologica}, 1 Q 96}
\end{quotex}


Inasmuch as Man shares in the nature of animals, Man is involved in the Lesser Struggle against material forces and enemies. \textbf{Work, temptation, and struggle are permanent aspects of life}.

Since Traditional political arrangements deal with reality, not speculative thought, the horizontal struggle involves the battle to protect and defend man's natural life. This includes the family, clan, tribe, and, of course, the City.

For a background on Guenon and Neoplatonism, see There is no God but God\footnote{\url{https://www.gornahoor.net/?p=3344}}


\flrightit{Posted on 2019-11-05 by Cologero}

\begin{center}* * *\end{center}

\begin{footnotesize}
\begin{sffamily}

\texttt{Ilo Stabet on 2021-03-16 at 05:04 said:}

The latter part (including the Aquinas quote) seems to contradict the Patristic notion (and even Scripture itself) in many regards, both before the Fall and before the Flood.

That before the Fall there was no death and therefore no struggle (and also no work, which is a form of struggle, as such -- hence the meaning of 'tilling the soil' after the Fall). That outside of paradise even the beasts we now know to be carnivores would eat only plants and fruits. And only after the flood was man (and beasts) allowed to eat flesh.

This struggle can also be seen as a result of disintegration between Heaven and Earth, and man's inability to mediate between the two, now requiring work, and technology, and so forth.

Or am I missing something?

\hfill

\texttt{Cologero on 2021-03-16 at 16:53 said:}

Yes, you are missing something.

The Lord God took the man and put him in the Garden of Eden to \textbf{work it and take care of it}. Genesis 2:15

\emph{So there was work}, although ``My yoke is easy, and my burden is light'' (Matthew 11:30).

then we read, The Lord God made garments of skin for Adam and his wife and clothed them. Genesis 3:21

Garments of skin come from dead animals. They weren't polyester.

Thomas' logic is impeccable. A lion that eats bananas and nuts is not a lion.

Animals cannot receive supernatural graces and their souls are not immortal.

Adam's soul is naturally immortal, but his body was not. Only supernatural grace can make a body immortal. But in Adam's case, I guess we will never know.

\hfill

\end{sffamily}
\end{footnotesize}

